\section{Financial statement modelling}

\begin{qfmethod}[A sales-based model with explicit financing and closing balances]\label{meth:integrated-pro-forma-statements}
\qffield{Statement} Forecast sales and operating drivers, then working capital, fixed assets, financing, taxes, cash and equity. Cash comes from the cash-flow statement; borrowing is an explicitly constrained financing decision. A model must identify infeasibility if available finance cannot fund the cash floor, rather than inserting unexplained debt, cash or retained earnings to force balance. Growth and required reinvestment must be consistent \citep[chap.~12, solutions~1--2, PDF p.~221]{statement-analysis-workbook-2020}; \citep[chap.~11, pp.~301--302]{investment-valuation-2012}.
\qffield{Practice} Start with Meridian's reconciled 2026 closing balances. Forecast 2027 sales at 1,100, COGS at 60 percent, cash-type operating expense at 18 percent, scheduled depreciation 55 and capex 90. Closing AR is 15 percent of sales, inventory one-third of COGS, AP 20 percent of COGS, and operating accruals 25 percent of operating expense. The 8 percent interest rate applies to opening debt 220; repayment of 20 and any revolver draw occur at year-end. Dividends are 45, tax is 25 percent of positive PBT and is paid currently. No deferred tax, OCI, equity issuance or other accounts are assumed.
\begin{center}\begin{tabular}{lrr|lrr}
Income/cash flow & 2026 & 2027 forecast & Balance sheet & Opening & Closing\\\hline
Sales & 1,000 & 1,100 & Cash & 87.5 & 96.05\\
COGS & 600 & 660 & AR & 150 & 165\\
Operating expense & 180 & 198 & Inventory & 200 & 220\\
Depreciation & 50 & 55 & PP\&E & 430 & 465\\
EBIT & 170 & 187 & Assets & 867.5 & 946.05\\
Interest & 20 & 17.6 & AP & 120 & 132\\
Tax & 37.5 & 42.35 & Accruals & 45 & 49.5\\
Net income & 112.5 & 127.05 & Debt & 220 & 200\\
CFO & 137.5 & 163.55 & Equity & 482.5 & 564.55\\
CFI & $-80$ & $-90$ & Contributed capital & 200 & 200\\
CFF & $-70$ & $-65$ & Retained earnings & 282.5 & 364.55
\end{tabular}\end{center}
\qffield{Reconciliation} CFO is $127.05+55-15-20+12+4.5=163.55$. Independently, receipts are 1,085, supplier payments 668, expense payments 193.5, interest 17.6 and tax 42.35. Cash closes at $87.5+163.55-90-20-45=96.05$; PP\&E at $430+90-55=465$; equity at $482.5+127.05-45=564.55$. Both sides total 946.05. With a 50 cash floor, the draw rule $\max(0,50-cash_{before\ draw})$ requests zero. A positive draw must fit a specified facility and feeds next year's interest; a facility limit is not an automatic source of unlimited money. The stress case below demonstrates that binding limit.
\end{qfmethod}

\begin{qfmethod}[Analyst forecast biases and practical remedies]\label{meth:forecast-bias-remedies}
\qffield{Statement} A numerical model does not remove subjective input choices. Anchoring on last year's result or management guidance, overconfidence in narrow ranges, confirmation bias, availability of recent events and pressure to agree with consensus can distort forecasts. Damodaran warns against choosing the answer before valuation and against treating consensus growth as independent evidence when analysts share information \citep[chap.~1, pp.~2--5; chap.~11, pp.~282--285]{investment-valuation-2012}.
\qffield{Practice} A small team has two hours to refresh Meridian's quarterly forecast. Management guides to 20 percent sales growth; the existing model assumes 10 percent and industry orders imply 5 percent. An analyst who keeps 20 percent solely to maintain access exhibits incentive/confirmation risk; mechanically averaging all three is not a remedy if they use the same order book. Recommend a driver bridge that separately records price, units, cancellations and new capacity, with independent evidence for each revision. Freeze a pre-results forecast and compare actual errors by driver, not only the total EPS miss.
\qffield{Alternatives and controls} Blindly using consensus is fast but inherits shared assumptions; a full new model each quarter is costly and may add false precision. Retain the reconciled model, independently challenge the largest cash and margin drivers, and report the base and adverse cases. Ask a second reviewer to identify evidence that would falsify the growth view, assign explicit ranges rather than implausibly precise point estimates, and update when that evidence arrives. Compare the forecast with mature peer margins and capacity constraints before accepting persistent excess returns. These controls reduce opportunities for bias; they do not claim to eliminate it or to guarantee forecast accuracy.
\end{qfmethod}

\begin{qfmethod}[Competition, price inflation and deflation in sales and cost forecasts]\label{meth:competition-price-cost-forecast}
\qffield{Statement} Translate competitive analysis into prices, volumes, unit costs and investment. Buyer power can limit price increases; supplier power raises input costs or tightens payment terms; substitutes cap feasible prices; entry adds capacity; rivalry can force discounts or selling expense. These are Porter's five forces, not five separate percentages to add to a growth forecast \citep[The Five Forces, five force descriptions]{competitive-forces-framework}. Damodaran likewise requires qualitative advantages to appear in margins, turnover, reinvestment or the duration of excess returns \citep[chap.~11, pp.~301--302]{investment-valuation-2012}.
\qffield{Practice} In a separate Meridian sensitivity, unit volume grows 4 percent and selling prices 3 percent: sales become $1{,}000(1.04)(1.03)=1{,}071.2$, not 1,070. If the same units' input costs rise 6 percent, COGS from a 600 base becomes $600(1.04)(1.06)=661.44$. Gross margin falls from 40 to 38.25 percent despite nominal sales growth. A powerful customer preventing any price rise gives revenue 1,040 and margin 36.4 percent. A substitute or new entrant makes passing through input inflation less credible; long-term supplier contracts may delay, rather than eliminate, the cost shock.
\qffield{Deflation and cash} With volume unchanged, prices falling 5 percent and unit costs falling only 2 percent give sales 950, COGS 588 and gross margin 38.11 percent. Historical-cost inventory can temporarily obscure that margin, while obsolescence may trigger NRV write-downs. Nominal receivables and debt do not fall automatically with prices. Forecast contract reset dates, inventory layers, wages and working-capital quantities separately; keep nominal revenue/cost forecasts consistent with nominal discount rates. Industry demand growth and company market-share changes require separate assumptions to avoid counting the same expansion twice.
\end{qfmethod}

\begin{qfmethod}[Choosing the explicit horizon and long-term forecast assumptions]\label{meth:forecast-horizon-terminal-assumptions}
\qffield{Statement} An explicit forecast horizon should cover the material transition to a sustainable operating state: product cycles, competitive advantage, capacity additions, contract expiries and financing changes matter more than a conventional five-year template. A cyclical peak is a poor terminal base. Damodaran links high-growth duration to firm size, current returns and sustainable competitive advantage, then requires stable-growth risk, reinvestment and financing assumptions to agree \citep[chap.~12, pp.~306--314]{investment-valuation-2012}.
\qffield{Practice} For the illustrative mature distributor, use 2027--2031 as the explicit planning horizon: 2027 has the detailed linked statements above; subsequent annual models should fade sales growth from 10 toward 3 percent and operating margin from 17 toward a supported 15 percent as competition absorbs excess returns. Those later years are a chosen modeling specification, not unexplained completed balance sheets. Use a longer horizon if a major plant or contract renewal prevents stabilization by 2031, or a shorter one if evidence supports an already stable business. Compare a three-year and seven-year transition as a sensitivity rather than letting the terminal date dictate the answer.
\qffield{Long-term mechanics} A stable nominal growth assumption must fit the currency and economy and remain below the discount rate. It also requires reinvestment: with next-period NOPAT 100, sustainable return on capital 10 percent and growth 3 percent, reinvestment is $100(0.03/0.10)=30$, leaving FCFF 70 and terminal value $70/(0.08-0.03)=1{,}400$ at 8 percent cost of capital. Raising growth to 4 percent requires reinvestment 40 and FCFF 60, giving 1,500, not 1,750 from holding cash flow fixed. If return on capital equals the 8 percent capital cost, terminal value for a fixed next-period NOPAT is $100/0.08=1{,}250$ regardless of sustainable growth below 8 percent: growth without excess returns creates no extra value. Neither an exit multiple nor a perpetual-growth formula excuses missing investment needs.
\end{qfmethod}
